\documentclass[10pt,a4paper]{amsart}
\usepackage{amsmath,amssymb,amsthm,graphicx}
% For commutative diagrams we use Xy-pic
\usepackage[all]{xy}

% We use 2cell for 2-commutative diagrams.
\xyoption{2cell}
% Theorem environments.
%
\theoremstyle{plain}
\newtheorem{theorem}[subsection]{Theorem}
\newtheorem{proposition}[subsection]{Proposition}
\newtheorem{lemma}[subsection]{Lemma}

\theoremstyle{definition}
\newtheorem{definition}[subsection]{Definition}
\newtheorem{example}[subsection]{Example}
\newtheorem{exercise}[subsection]{Exercise}
\newtheorem{situation}[subsection]{Situation}

\theoremstyle{remark}
\newtheorem{remark}[subsection]{Remark}
\newtheorem{remarks}[subsection]{Remarks}

\numberwithin{equation}{subsection}

% Macros
%
\def\lim{\mathop{\mathrm{lim}}\nolimits}
\def\colim{\mathop{\mathrm{colim}}\nolimits}
\def\Spec{\mathop{\mathrm{Spec}}}
\def\Hom{\mathop{\mathrm{Hom}}\nolimits}
\def\Ext{\mathop{\mathrm{Ext}}\nolimits}
\def\SheafHom{\mathop{\mathcal{H}\!\mathit{om}}\nolimits}
\def\SheafExt{\mathop{\mathcal{E}\!\mathit{xt}}\nolimits}
\def\Sch{\mathit{Sch}}
\def\Mor{\mathop{\mathrm{Mor}}\nolimits}
\def\Ob{\mathop{\mathrm{Ob}}\nolimits}
\def\Sh{\mathop{\mathit{Sh}}\nolimits}
\def\NL{\mathop{N\!L}\nolimits}
\def\CH{\mathop{\mathrm{CH}}\nolimits}
\def\proetale{{pro\text{-}\acute{e}tale}}
\def\etale{{\acute{e}tale}}
\def\QCoh{\mathit{QCoh}}
\def\Ker{\mathop{\mathrm{Ker}}}
\def\Im{\mathop{\mathrm{Im}}}
\def\Coker{\mathop{\mathrm{Coker}}}
\def\Coim{\mathop{\mathrm{Coim}}}

% Boxtimes
%
\DeclareMathSymbol{\boxtimes}{\mathbin}{AMSa}{"02}

%
% Macros for moduli stacks/spaces
%
\def\QCohstack{\mathcal{QC}\!\mathit{oh}}
\def\Cohstack{\mathcal{C}\!\mathit{oh}}
\def\Spacesstack{\mathcal{S}\!\mathit{paces}}
\def\Quotfunctor{\mathrm{Quot}}
\def\Hilbfunctor{\mathrm{Hilb}}
\def\Curvesstack{\mathcal{C}\!\mathit{urves}}
\def\Polarizedstack{\mathcal{P}\!\mathit{olarized}}
\def\Complexesstack{\mathcal{C}\!\mathit{omplexes}}
% \Pic is the operator that assigns to X its picard group, usage \Pic(X)
% \Picardstack_{X/B} denotes the Picard stack of X over B
% \Picardfunctor_{X/B} denotes the Picard functor of X over B
\def\Pic{\mathop{\mathrm{Pic}}\nolimits}
\def\Picardstack{\mathcal{P}\!\mathit{ic}}
\def\Picardfunctor{\mathrm{Pic}}
\def\Deformationcategory{\mathcal{D}\!\mathit{ef}}
\newenvironment{reference}{}{}
\newenvironment{slogan}{}{}
\newenvironment{history}{}{}
\usepackage{cleveref}
\crefname{originallemma}{original lemma}{original lemmas}
\Crefname{originallemma}{Original lemma}{Original lemmas}
\makeatletter
\newlabel{lemma-recurring-real}{{0C5V}{source}{}{sourcepointer-1}{}}
\newlabel{lemma-recurring-real@cref}{{[originallemma][0][]0C5V}{[1][0][]source}{}{}{}}
\makeatother
\crefname{originallemma}{original lemma}{original lemmas}
\Crefname{originallemma}{Original lemma}{Original lemmas}
\makeatletter
\newlabel{lemma-recurring-symmetric-real}{{0C5X}{source}{}{sourcepointer-2}{}}
\newlabel{lemma-recurring-symmetric-real@cref}{{[originallemma][0][]0C5X}{[1][0][]source}{}{}{}}
\makeatother
\title{The Picard group of a numerical type is a finitely generated abelian group of rank one}
\author{The Stacks Project Authors}
\date{stacks-project:models.tex:local-snapshot-20261012; GNU FDL 1.2 or later}
\begin{document}
\maketitle
\noindent\textit{Editorial note.} The counted unit is the Stacks Project Lemma lemma-picard-rank-1 at source lines 1128-1132 together with its proof at source lines 1134-1141, a straight adjacent pairing in the chapter file. The chapter is the upstream file models.tex; the pool directory carries the longer name stacks-models and is not used as the chapter name anywhere in this card.\par\medskip
\noindent\textit{Editorial note.} One context span carries the definition of the Picard group of a numerical type at lines 1111-1126, which the retained statement and proof both refer to; the span is the complete definition environment from its opening to its closing marker, it is context rather than a counted unit, and it does not overlap the counted interval. The first revision had pointed the span at lines 1098-1103, which is in fact a remark about equality of genera, and the engine refused the file because that span began inside a remark environment and left it unclosed; this revision carries the correct span.\par\medskip
\noindent\textit{Editorial note.} The retained proof names two lemmas of the same chapter, the real and the symmetric forms of the recurrence statement about the rank of the intersection matrix, both counted by separate cards of this lane; they are declared as non-load-bearing navigation pointers with their stable Stacks tags. The definition retained as context is deliberately not declared as a pointer as well, so that no label is defined twice.\par\medskip
\noindent\textit{Editorial note.} No citation occurs in the retained spans, so no bibliography entry is supplied.\par\medskip
\noindent\textit{Allowed standard prerequisites (editorial).}\begin{itemize}
\item Numerical types, their recurrence conditions and their Picard group. A numerical type is a collection of integers n, m of subscript i, a of subscript i j, w of subscript i and g of subscript i whose intersection numbers are symmetric, whose diagonal entries are nonpositive with vanishing row sums weighted by the multiplicities, and whose weights divide the corresponding intersection numbers; its Picard group is the cokernel of the matrix of the divided intersection numbers. The chapter proves separately that under the recurrence conditions the intersection matrix has rank one less than its size, in two versions, one for real and one for symmetric data. The author's own definitions and those two lemmas are the allowed inputs.. The counted lemma computes the rank of the Picard group of a numerical type, which the section needs when it compares Picard groups of types related by a contraction and when it identifies the Picard group in the case of nonpositive genus.
\end{itemize}\medskip
\ifdefined\hypertarget\hypertarget{sourcepointer-1}{}\fi
\noindent\textit{Source reference.} Original lemma 0C5V: focused-seeds/AGP/sources/stacks-snapshot-20261011/tags.txt, source lines 15275--15275. Navigation pointer invoked by the retained proof for the rank of the intersection matrix under the real form of the recurrence conditions; that lemma is counted by a separate card of this lane, so its proof is not reproduced here. The number is the stable Stacks tag of the target, looked up in the shared tag index.\par\medskip
\ifdefined\hypertarget\hypertarget{sourcepointer-2}{}\fi
\noindent\textit{Source reference.} Original lemma 0C5X: focused-seeds/AGP/sources/stacks-snapshot-20261011/tags.txt, source lines 15277--15277. Navigation pointer invoked by the retained proof for the rank of the intersection matrix under the symmetric form of the recurrence conditions; that lemma is counted by a separate card of this lane, so its proof is not reproduced here. The number is the stable Stacks tag of the target, looked up in the shared tag index.\par\medskip
\begin{definition}
\label{definition-picard-group}
Let $n, m_i, a_{ij}, w_i, g_i$ be a numerical type $T$. The
{\it Picard group of $T$} is the cokernel of the matrix
$(a_{ij}/w_i)$, more precisely
$$
\Pic(T) =
\Coker\left(
\mathbf{Z}^{\oplus n} \to \mathbf{Z}^{\oplus n},\quad
e_i
\mapsto
\sum \frac{a_{ij}}{w_j}e_j
\right)
$$
where $e_i$ denotes the $i$th standard basis vector for $\mathbf{Z}^{\oplus n}$.
\end{definition}

\begin{lemma}
\label{lemma-picard-rank-1}
Let $n, m_i, a_{ij}, w_i, g_i$ be a numerical type $T$.
The Picard group of $T$ is a finitely generated abelian group of rank $1$.
\end{lemma}

\begin{proof}
If $n = 1$, then $A = (a_{ij})$ is the zero matrix and
the result is clear. For $n > 1$ the matrix $A$ has rank
$n - 1$ by either Lemma \ref{lemma-recurring-real} or
Lemma \ref{lemma-recurring-symmetric-real}.
Of course the rank is not affected by scaling the rows
by $1/w_i$. This proves the lemma.
\end{proof}

\end{document}
